% =========================
% Chapter 5: Rule Replay, Voting System, and Counterfactuals
% File: 5_rule_reply_volting_system_and_counterfactuials.tex
% =========================

% ---- Chapter opening positioning paragraph (results-focused, protocol explicit)
Following the official rule definitions, we perform season--week counterfactual replays under (i) rank-based aggregation and (ii) percent-based aggregation. We initialize each replay with the observed eligibility labels, but update the active set week-by-week according to the counterfactual eliminations under the chosen rule. When a counterfactual path keeps an observed-exit contestant active beyond the observed data horizon, we apply a conservative, minimal-information imputation solely to complete the replay (not to forecast future performance).
As a reproducible sanity-check baseline (BL-0), we use judge-only information with a uniform fan proxy
and report elimination hit-rate (KPI1) and the rule divergence rate (KPI3/FlipRate).
This chapter summarizes cross-season differences (Section~5.1), explains divergences via controversial case studies (Section~5.2),
evaluates the incremental impact of bottom-two and judges' save (Section~5.3), and synthesizes implications (Section~5.4).

\subsection{Rank vs. Percent Across Seasons}\label{subsec:rank_vs_percent_across_seasons}
Results in this section are first reported under BL-0 as a reproducible sanity check; Task A inferred votes are used in subsequent replays.
% 这里放：active set 口径、single-elim eligibility、special exits、tie-handling、以及“同一输入不同规则”的对比原则。
% 目标：评审只看这一段就不会质疑你们分母与口径。
  
% ---------------- FIG 5.1: FlipRate by season (eligible weeks only) ----------------
% 插入点 A：在 Cross-season summary metrics (KPI1 & KPI3) 第一段之后
% 用途：先给读者“全局敏感度”直观印象
% Fig 5.1 fig_fliprate_by_season.pdf（FlipRate by season，eligible only）
% -------------------------------------------------------------------------------
% Insert here to support KPI3/FlipRate cross-season comparison.
\begin{figure}[H]
    \centering
    \includegraphics[width=0.95\linewidth]{output/fig/fig5\_1\_fliprate\_by\_season.png}
    \caption{FlipRate by season under comparable (eligible) weeks.}
    \label{fig:fliprate_by_season}
\end{figure}
% -------------------------------------------------------------------------------

\subsubsection{Replay protocol and eligibility}\label{subsubsec:replay_protocol_and_eligibility}

\paragraph{Eligibility and special episodes.}
We treat eligibility as part of the replay interface: all week-level comparisons and KPI denominators are defined over the active set,
and special episodes (no-elimination weeks, multi-elimination weeks, withdrawals, and finals) are handled explicitly rather than being forced into a single-elimination template.
This design ensures that rule comparisons reflect the mechanics actually used on the show.

\paragraph{Aggregation rules (rank vs. percent).}
We replay eliminations under two official aggregation paradigms: rank-based and percent-based. For percent-based aggregation,
the combined score is
\[
S^{(P)}_{i,t}(\alpha)=\alpha q^{J}_{i,t}+(1-\alpha)\hat p_{i,t},
\]
and the replay eliminates the $m_{s,t}$ contestants with smallest $S^{(P)}_{i,t}(\alpha)$, where $m_{s,t}$ is the number of eliminations in week $(s,t)$.
For rank-based aggregation, we replace the score shares by within-week ranks and apply the same elimination logic.
Across both paradigms, ties are broken deterministically using a fixed secondary key (e.g., judges' score, then stable contestant ID), to preserve reproducibility.

\subsubsection{Cross-season summary metrics}\label{subsubsec:cross_season_summary_metrics}
We compute replay-based summary metrics over counted (season, week) instances under a unified eligibility convention.
KPI1 (Jaccard set agreement) summarizes week-level agreement between the observed eliminated set and the replayed eliminated set,
including no-elimination weeks via an explicit empty-set convention.
KPI3 (FlipRate) measures how often the predicted eliminated set differs between percent-based and rank-based replays under the same inputs.

\paragraph{Interpretation.}
We measure the bias to determine the degree of its inclination towards the audience, where a higher bias value indicates a greater tendency to favor the audience. 
According to Figure~\ref{fig:bias_index_comparison_four_methods_alpha0p5} we find that the \textit{percentage} method is more inclined towards the audience, yet the actual difference between the two methods is negligible compared with the \textit{PB2} we designed in \ref{subsubsec:d2_ruleB}.



\subsection{Controversial Contestants: Mechanism-driven Case Studies}\label{subsec:controversial_contestants_case_studies}

Aggregate metrics such as FlipRate identify \emph{where} rules diverge, but they do not explain \emph{why} a particular outcome
was perceived as controversial. We therefore adopt mechanism-driven case studies: each selected contestant represents a distinct
interaction between judges' performance signals and audience support under the incumbent ``combined Bottom-2 + judges' save'' pipeline
and under alternative aggregation semantics (rank vs.\ percent).

\paragraph{Selection rationale (what makes a case ``controversial'' in our framework).}
A contestant is considered controversial if their survival path is supported by a persistent mismatch between judges' signals and inferred fan support,
so that small changes in aggregation (rank vs.\ percent) or stage-ordering (when judges intervene) can plausibly alter the elimination boundary.
Operationally, we focus on seasons/episodes with elevated FlipRate and on contestants frequently cited in public discourse about judge--fan disagreement.
We analyze four representative contestants: Jerry Rice (Season~2), Billy Ray Cyrus (Season~4), Bristol Palin (Season~11), and Bobby Bones (Season~27).
All summaries are computed from \texttt{clean\_long\_data\_new1.csv} using \texttt{season}, \texttt{week}, \texttt{celebrity\_name}, \texttt{placement},
\texttt{elim\_week}, and judges-derived variables (\texttt{J\_total}, \texttt{J\_pct}, \texttt{active}).

\paragraph{Protocol (kept fixed across cases).}
For each contestant we report: (i) their within-season judges trajectory (e.g., week-level position under $q^J_{i,t}$), (ii) their inferred fan-support trajectory
through $\hat p_{i,t}$ (Task~A), and (iii) counterfactual replays under percent vs.\ rank aggregation with and without the judges' save module.
The objective is not to ``explain away'' controversy, but to isolate the minimal mechanism that makes the outcome boundary-sensitive.

\paragraph{Re-rank outcomes between \textit{percentage} and \textit{rank} rules for controversial contestants.} 
When we re-rank Jerry Rice under the \textit{percentage} rule and re-rank Billy Ray Cyrus, Bristol Palin and Bobby Bones under the \textit{rank} rule, 
we find that the final rankings of all four contestants drop significantly in the Tabel~\ref{tab:rank_percent}. This indicates a notable discrepancy in the tolerance for controversial 
contestants after the rule adjustment, and we will analyze these characteristics in the subsequent sections.

\begin{table}[htbp]
  \centering
  \caption{Rank and Percent Data Table} % 表格标题，可自行修改
  \begin{tabular}{lcccccc} % l=左对齐，c=居中对齐，7列对应7个字段
    \toprule
    name & season & canon & rank & percent & percent\_last\_two & bottom\_two \\
    \midrule
    Billy Ray Cyrus & 4 & 5 & 10 & 5 & 10 & 6 \\
    Bobby Bones    & 27 & 1 & 4 & 4 & 1 & 5 \\
    Bristol Palin  & 11 & 3 & 3 & 3 & 3 & 4 \\
    Bristol Palin  & 15 & 9 & 13 & 9 & 13 & 12 \\
    Jerry Rice     & 2 & 2 & 2 & 3 & 3 & 4 \\
    \bottomrule
  \end{tabular}
  \label{tab:rank_percent} % 表格交叉引用标签，可自行修改
\end{table}

\paragraph{Case 1: Jerry Rice (Season~2) --- early-era aggregation sensitivity.}
This case illustrates how early-season aggregation semantics can amplify judge--fan discordance: a contestant can remain viable when fan support offsets weak
judges' scores, yet the elimination boundary becomes sensitive to whether scores are combined as ranks or as normalized shares.
In our replay framework, small changes in aggregation can change which contestants enter the Bottom-2 risk set, which then cascades through subsequent weeks.

\paragraph{Case 2: Billy Ray Cyrus (Season~4) --- boundary effects without outcome reversal.}
Not every ``controversial'' public narrative implies a counterfactual outcome reversal. This case serves as a negative control: despite persistent discussion,
our replays indicate that the elimination week can remain stable under the incumbent pipeline.
We retain it to show that controversy perception can arise from judges--fan disagreement even when the realized elimination is not easily flipped by reasonable rule perturbations.

\paragraph{Case 3: Bristol Palin (Season~11) --- popularity protection under fan-sensitive boundaries.}
This case represents a common mechanism in celebrity competitions: strong fan support can keep a weak performer out of the Bottom-2 set for multiple weeks,
especially when the risk set is defined by a combined score.
When a judges' save is applied only after the Bottom-2 is formed, the mechanism protects technical merit only at the final step and cannot prevent fan support from shaping entry into the risk set.

\paragraph{Case 4: Bobby Bones (Season~27) --- late-era mismatch and institutional interpretation.}
This case highlights that the same observed trajectory can be interpreted differently depending on institutional design.
Under the incumbent ``combined Bottom-2 + judges' save'' pipeline, fan support can strongly affect who enters the risk set;
under a professionalism-first gate (Task~D2), the risk set is defined by judges alone, which changes whether fan support can ``rescue'' a weak performer.
We use this case to motivate why stage-ordering (gate vs.\ final save) is a central institutional lever rather than a cosmetic change.



\paragraph{Mechanisms synthesized across cases.}
Across the four examples we observe three recurring mechanisms:
(i) \textbf{aggregation semantics} (rank vs.\ percent) can shift the boundary of who is at risk;
(ii) \textbf{risk-set formation} is often more consequential than the final tie-break or save decision; and
(iii) \textbf{stage-ordering} determines whether judges' influence acts as a last-step safeguard (incumbent) or as a professionalism-first gate (Task~D2).
These mechanisms explain why controversies cluster around boundary-sensitive weeks and motivate the institutional redesign evaluated in Section~7.


\subsection{Bottom-two and Judges' Save Mechanism}\label{subsec:bottom_two_and_judges_save}

Some later seasons introduce a two-stage mechanism: (i) identify a bottom subset using the vote aggregation rule, then (ii) select the eliminated contestant from that subset via judges' save.
To isolate the incremental effect of this change, we implement a modular judges' save layer that can be toggled on top of either percent-based or rank-based aggregation.

\textbf{Module definition.} In a week with a bottom-two format, the replay first forms a bottom-two set $B_{s,t}$ using the chosen aggregation rule.
The saved contestant is the one with higher judges' score $J_{i,t}$ (equivalently higher $q^J_{i,t}$), and the other is eliminated.
This construction cleanly separates \emph{how the bottom set is formed} from \emph{how the final elimination is chosen}, enabling controlled counterfactual comparisons.

\textbf{Evaluation questions.} We quantify (a) how often the judges' save reverses the elimination implied by votes alone, (b) whether the save disproportionately protects contestants with strong judges' scores but weak fan support, and (c) how the mechanism affects overall rule divergence and late-season stability.
These results connect directly to management decisions about perceived fairness versus entertainment value.\cite{brandt2016comsoc,saari2001decisions}

\textbf{Illustrative case-study replay.} As a directional example, we run a Bottom-2 + judges' save replay with $\alpha=0.5$
(equal weight between $q^J_{i,t}$ and $\hat p_{i,t}$ in the combined score definition above) on Seasons~2/4/11/27.
Under this module, contestants whose survival historically relies on fan support tend to be eliminated earlier: Jerry Rice (S2) is eliminated in Week~7,
Bristol Palin (S11) in Week~9, and Bobby Bones (S27) in Week~8 in our replay, while Billy Ray Cyrus (S4) remains unchanged (Week~8).

\textbf{Caveat.} These outcomes can change for marginal weeks under different $\alpha$ values (we chose $\alpha \in [0.3, 0.7]$ in practice), tie-breaking conventions, and the semantics of
multi-elimination weeks (we apply eliminations sequentially when multiple eliminations occur), so we treat the above as illustrative rather than universal.




% ==========================规则概述与反事实分析=============================
\subsection{Synthesis of Findings from Rules}\label{subsec:synthesis_of_findings_from_rules}

We synthesize the rule replay evidence into actionable takeaways for both mechanism understanding and rule design.
First, rank versus percent aggregation differences are not uniformly distributed across seasons; they are amplified in weeks where contestants cluster near the elimination boundary, making the system sensitive to small score-share changes.
Second, the bottom-two + judges' save module acts as a structural preference for judges at the final decision margin, reducing the extent to which fan support can rescue a contestant once the contestant enters the bottom set.
Third, these rule interactions motivate the professionalism-first institutional redesign in Task D-2: moving the judges' influence from a final-step safeguard to a gating mechanism, while preserving watchability through structured late-stage suspense mechanisms evaluated via replay and stress testing.

% ---------------- OPTIONAL: Summary box (if you want a visual takeaway) ----------------
% Consider a small boxed summary here for Memo-ready bullets.
% -------------------------------------------------------------------------------




% =========================下面是我们粘图片的位置注释=================

% % ---------------- TAB 5.1: KPI1 hit-rate (rank vs percent) with era slices ----------------
% 用途：给出 KPI1（hit-rate）对照表，体现“fidelity vs sensitivity”
% Tab 5.1 tab_baseline_consistency.tex（Overall + era slices + #eligible weeks）
% -------------------------------------------------------------------------------
% % Insert here to report BL-0 sanity-check baseline consistency (eligible weeks only).
% \begin{table}[H]
%     \centering
%     %\begin{tabular}{lrrrr}
\hline
Slice & $N_{\mathrm{weeks}}$ & Hit-Rate (Rank) & Hit-Rate (Percent) & FlipRate \\
\hline
Overall & 264 & 0.364 & 0.375 & 0.038 \\
Rank-era (S1-2,28-34) & 66 & 0.348 & 0.364 & 0.015 \\
Percent-era (S3-27) & 198 & 0.369 & 0.379 & 0.045 \\
\hline
\end{tabular}
%     \caption{Elimination hit-rate (KPI1) for rank-based vs.\ percent-based replays (eligible weeks only), including era slices.}
%     \label{tab:kpi1_hitrate_era}
% \end{table}
% % -------------------------------------------------------------------------------

% % ---------------- FIG 5.2: Era cutoff stress test (27/28/29) ----------------
% 插入点 C：在 Era cutoff assumption and stress-test plan 段落之后
% 用途：一眼让评审看到“cutoff 变动对结论影响多大”，避免被质疑
% Fig 5.2 fig_cutoff_stress_test.pdf（cutoff=27/28/29 下的 FlipRate 或 hit-rate）
% -------------------------------------------------------------------------------
% % Insert here to show robustness of KPI1/KPI3 to era boundary assumption.
% \begin{figure}[H]
%     \centering
%     %\includegraphics[width=0.80\linewidth]{figures/fig_cutoff_stress_test.pdf}
%     \fbox{\parbox{0.80\linewidth}{\centering Placeholder: Fig 5.2 Era Cutoff Stress Test}}
%     \caption{Sensitivity of replay metrics to the assumed era cutoff (e.g., 27/28/29). Error bars denote bootstrap uncertainty over weeks.}
%     \label{fig:cutoff_stress_test}
% \end{figure}
% % -------------------------------------------------------------------------------

% % ---------------- FIG 5.3: FlipRate by season phase (early/mid/late) ----------------
% （P1 可选）插入点 D：如果你们写了“late-season 更敏感”的句子，就在那句之后插
% ---------------·------------------------------------------------
% % Only include if the text claims phase dependence (late-season sensitivity).
% \begin{figure}[H]
%     \centering
%     %\includegraphics[width=0.80\linewidth]{figures/fig_fliprate_by_phase.pdf}
%     \fbox{\parbox{0.80\linewidth}{\centering Placeholder: Fig 5.3 FlipRate by Phase}}
%     \caption{FlipRate by normalized season phase (eligible weeks only).}
%     \label{fig:fliprate_by_phase}
% \end{figure}
% % -------------------------------------------------------------------------------

% % ---------------- FIG 5.4x: Case timeline (judges vs fan) ----------------
% % 5.2 Controversial contestants case studies — 插图/表注释占位
% % 每个案例块建议插两张图（你们先留空位，等跑出来填）
% % 插入点 E：案例叙事中“Observed trajectory”之后
% % Fig 5.4a / 5.4b fig_case_<name>_timeline.pdf（judges rank、fan rank、淘汰标记）
% % -------------------------------------------------------------------------------
% % Insert inside each case study after "Observed trajectory".
% \begin{figure}[H]
%     \centering
%     %\includegraphics[width=0.90\linewidth]{figures/fig_case_JerryRice_timeline.pdf}
%     \fbox{\parbox{0.90\linewidth}{\centering Placeholder: Fig 5.4x Case Timeline}}
%     \caption{Case timeline example: week-by-week judges rank and inferred fan rank, with elimination markers.}
%     \label{fig:case_timeline_placeholder}
% \end{figure}
% % -------------------------------------------------------------------------------


% % ---------------- TAB 5.2x: Case counterfactual summary table ----------------
% % 插入点 F：在“Counterfactual replays”段落之后
% % Tab 5.2（小表） tab_case_<name>_counterfactual.tex：列出“真实淘汰 vs percent replay vs rank replay（含 save/不含 save）”
% % -------------------------------------------------------------------------------
% % Insert after "Counterfactual replays" to make the narrative audit-able.
% \begin{table}[H]
%     \centering
%     %\input{tables/tab_case_JerryRice_counterfactual.tex}
%     \caption{Case counterfactual summary: observed vs.\ percent vs.\ rank (and save variants if applicable).}
%     \label{tab:case_counterfactual_placeholder}
% \end{table}
% % -------------------------------------------------------------------------------


% % ---------------- FIG 5.5: Judges' save effect (elimination changed rate) ----------------
% % 5.3 Bottom-two and Judges' Save Mechanism — 插图注释占位
% % 插入点 G：在 “Evaluation questions” 段落之后
% % Fig 5.5 fig_save_fliprate.pdf：save 导致淘汰改变的比例（按 season 或按 week）
% % ------------------------------------------------------------------------------
% % Insert here to quantify how often judges' save changes the eliminated contestant.
% \begin{figure}[H]
%     \centering
%     %\includegraphics[width=0.85\linewidth]{figures/fig_save_fliprate.pdf}
%     \fbox{\parbox{0.85\linewidth}{\centering Placeholder: Fig 5.5 Judges' Save Effect}}
%     \caption{Impact of judges' save: rate at which the eliminated contestant changes relative to vote-only elimination.}
%     \label{fig:judges_save_effect}
% \end{figure}
% % -------------------------------------------------------------------------------


% % ------------------------FIG 5.6 fig_save_protection_profile.pdf------------------------
% （P1 可选）如果你们要支撑“save 更偏向保护技术高的人”：再加
% Fig 5.6 fig_save_protection_profile.pdf：被 save 的人通常 judges rank 更高但 fan rank 更低的比例（箱线/散点）
% ------------------------------------------------------------------------------